\begin{abstract}
Large language models (LLMs) are now built into everyday tools, yet they
remain unreliable when an answer depends on information that is very large
or rarely seen in their training data, such as an organization's own
systems and the execution traces those systems produce.
Retrieval-augmented generation (RAG) helps by supplying relevant documents,
but documents alone cannot tell the model what actually happened in a
specific run.
We present \sysname, a framework that answers plain-language questions
about Linux kernel traces. Java analyzers built on Trace Compass process
the whole trace and compute measured findings, such as how each thread's
time divides between running, waiting for CPU and being blocked, and which
thread wakes which; RAG adds domain and program documentation; and an LLM
explains the findings.
We evaluate it on LTTng traces of a multithreaded program with one slow
thread, recorded at 2, 10 and 30 minutes, using 13 questions covering
comprehension, root cause, dependencies, quantitative analysis and
remediation.
With the same model (GPT-5.6), an LLM given raw trace excerpts scores 35\%
and an LLM with RAG scores 68\%. \sysname answers all 13 questions
correctly on every trace, with timings that match the measured ground
truth.
The results suggest that grounding an LLM in computed evidence, not only
retrieved text, lets developers without tracing expertise diagnose
performance problems by asking questions.
\end{abstract}

\begin{IEEEkeywords}
large language models, retrieval-augmented generation, execution traces,
LTTng, Trace Compass, performance analysis
\end{IEEEkeywords}
